\documentclass[submission,copyright,creativecommons]{eptcs}
\providecommand{\event}{PCC 2026: Proof, Computation, and Complexity}
\usepackage{underscore}

\title{The Adaptive Linguistic Security Framework (ALSF):\\ Deterministic Suppression of Large Language Model Hallucination via Axiomatic Grammars}
\author{Rafael D. De Paz
\institute{Vanguard Research Array \\ Sovereign Architecture Group \\ Las Vegas, NV}
\email{rafael@rdepaz.com}
}
\def\titlerunning{The Adaptive Linguistic Security Framework (ALSF)}
\def\authorrunning{R. D. De Paz}


\usepackage[top=1in, bottom=1in, left=1in, right=1in]{geometry}
\usepackage{draftwatermark}
\SetWatermarkText{SOVEREIGN PROTOCOL // ID: B96B98A}
\SetWatermarkScale{0.5}
\SetWatermarkLightness{0.94}
\begin{document}
\maketitle

\begin{abstract}
The fundamental fragility of continuous-space probabilistic models (LLMs) is their inherent capacity for "hallucination"—the generation of syntactically valid but semantically or factually incoherent outputs. This paper introduces the Adaptive Linguistic Security Framework (ALSF), a deterministic validation layer operating below the threshold of standard semantic generation. By imposing strict, computable L0-L1 axiomatic grammars as hard boundaries acting as a "Justiciary", we demonstrate that it is mathematically possible to guarantee 100\% hallucination suppression in targeted execution domains (e.g., Code Execution, API invocation). We provide formal proofs of the ALSF architectural sieve and discuss its implementation within The Logos Kernel Protocol infrastructure.
\end{abstract}

\end{document}
